Accurate multi-horizon load forecasting is the critical enabler of proactive radio resource allocation, providing the necessary lead time for APs to coordinate before crowds assemble. Table~\ref{tab:pred} and Figure~\ref{fig:pred} summarize the forecasting performance across all horizons ($\Delta=6$, spanning 30 minutes).

ST-GCAT achieves an overall multi-step RMSE of $\mAOneRmse$, outperforming the persistence baseline ($\mPersistRmse$, a $35\%$ error reduction), the temporal-only LSTM baseline ($\mLSTMDRLRmse$, $15\%$), and the schedule-unaware ST-GCAT model ($\mATwoRmse$, $31\%$); these values are those of the decoupled ($\eta=0$) row of Table~\ref{tab:pred}, and the predictor of the proposed configuration, which has the same architecture, reaches $\mPropRmse$. The predictive advantage of ST-GCAT is particularly pronounced during critical scheduled surge windows: ST-GCAT achieves a Surge-RMSE of $\mAOneSurgeRmse$, representing a $39\%$ reduction compared to persistence ($\mPersistSurgeRmse$), $24\%$ compared to the LSTM baseline ($\mLSTMDRLSurgeRmse$), and $33\%$ compared to the schedule-unaware model ($\mATwoSurgeRmse$). Outside of surge windows, the LSTM and the ST-GCAT variants that use the schedule have a Normal-RMSE of $0.111$--$0.118$, against $\mPersistNormalRmse$ for persistence and $\mATwoNormalRmse$ for ST-GCAT without schedule.

These results reveal two fundamental insights into spatio-temporal network forecasting:
\begin{enumerate}
  \item \textbf{The Essential Role of Schedule Cross-Attention}: Without schedule conditioning, the graph-attention architecture alone achieves an RMSE of $0.175$ and Surge-RMSE of $0.243$. This indicates that, in our simulator (where the timetable is exact), the cross-attention over upcoming timetable metadata is the main source of the gain, and that the graph and temporal layers alone do not anticipate announced crowd arrivals from historical trends.
  \item \textbf{Horizon-Adaptive Dynamics}: At the immediate one-step horizon ($h=1$, 5 minutes ahead), historical load autocorrelation dominates: the RMSE is $\mAOneRmseHOne$ for ST-GCAT against $\mPersistRmseHOne$ for persistence and $\mLSTMDRLRmseHOne$ for the LSTM, i.e.\ ST-GCAT is \emph{less} accurate at $h=1$. Its lower overall RMSE therefore comes from the longer horizons, where the timetable carries information that the history does not; we did not break the error down per horizon.
\end{enumerate}
